\section{Reshaping two intercepts}\label{sec:reshaping}

We now change two intercepts of a simplex while preserving its volume, the
cardinality of every exponent set, and the graded inclusion.  The replacement
will also preserve the implication from monomial jet surjectivity back to the
original spaces.  The important freedom is that the new small intercept can
be prescribed as a real number; the argument can therefore be iterated even
when the preceding step has produced irrational intercepts.

The node, diagonal compression, and primitive-direction choice adapt the
corresponding constructions of
\citet[Propositions~3.2 and~3.4 and Lemma~4.3]{OpenAI2026Surfaces}.
Here they are applied directly to polynomial spaces in two variables, with
the remaining exponents retained as labels.  We give the argument in full.

\begin{proposition}[Two-intercept reshaping]\label{prop:reshape}
Let $n\ge2$, let $A,B,a_3,\ldots,a_n$ be positive real numbers, and let
$(S_k)_{k\ge1}$ be finite exponent sets satisfying
\[
 S_k\subset k\Delta(A,B,a_3,\ldots,a_n)\cap\Z^n,
 \qquad S_k+S_l\subset S_{k+l}\quad(k,l\ge1).
\]
Set $H=AB$ and $d=\max(3/A,2/B)$.  For every real number $w>0$ with
$dw<1/16$, there are finite sets $(S'_k)_{k\ge1}$ such that
\begin{align}
 S'_k&\subset k\Delta(w,H/w,a_3,\ldots,a_n)\cap\Z^n,
                                                        \label{eq:reshape-bound}\\
 \#S'_k&=\#S_k,
 \qquad S'_k+S'_l\subset S'_{k+l}.                       \label{eq:reshape-graded}
\end{align}
Moreover, for every $k\ge1$, $m\ge0$, and $r\ge1$, if $\M(S'_k)$ surjects
onto jets through degree $m$ at some $r$ distinct points of $\Torus^n$,
then $\M(S_k)$ does so at some $r$ distinct points of $\Torus^n$.
\end{proposition}

\newcommand{\ReshapingNodeFigure}{%
\begin{figure}[H]
\centering
\begin{tikzpicture}[x=.88cm,y=2.8cm,font=\small,
  polyshape/.style={draw=teal!65!black,fill=teal!10,line width=.8pt},
  slice/.style={draw=orange!85!black,line width=1.2pt}]
  \begin{scope}
    \node at (1.4,.77) {$Q$};
    \draw[->,gray!65] (-.2,0)--(3.35,0);
    \draw[->,gray!65] (0,-.04)--(0,.63);
    \path[polyshape] (0,0)--(3,0)--(.5,.5)--(0,.375)--cycle;
    \node[below left] at (0,0) {$0$};
    \node[below] at (3,0) {$(a,0)$};
    \node[above right] at (.5,.5) {$(c,c)$};
    \node[left] at (0,.375) {$(0,b)$};
    \draw[slice] (0,.18)--(.26,.44);
    \node at (1.4,-.21) {diagonal slices};
  \end{scope}
  \begin{scope}[xshift=4.65cm]
    \node at (1.35,.77) {after the shear};
    \draw[->,gray!65] (-.7,0)--(3.35,0);
    \draw[->,gray!65] (0,-.04)--(0,.63);
    \path[polyshape] (0,0)--(3,0)--(0,.5)--(-.375,.375)--cycle;
    \node[below right] at (0,0) {$0$};
    \node[below] at (3,0) {$(a,0)$};
    \node[above right] at (0,.5) {$(0,c)$};
    \node[left] at (-.375,.375) {$(-b,b)$};
    \draw[slice] (-.18,.18)--(-.18,.44);
    \node at (1.35,-.21) {vertical slices};
  \end{scope}
  \begin{scope}[xshift=9.3cm]
    \node at (1.35,.77) {$P$};
    \draw[->,gray!65] (-.7,0)--(3.35,0);
    \draw[->,gray!65] (0,-.04)--(0,.63);
    \path[polyshape] (-.375,0)--(3,0)--(0,.5)--cycle;
    \node[below left] at (-.375,0) {$(-b,0)$};
    \node[below] at (3,0) {$(a,0)$};
    \node[above right] at (0,.5) {$(0,c)$};
    \draw[slice] (-.18,0)--(-.18,.26);
    \node at (1.35,-.21) {slices moved to height $0$};
  \end{scope}
  \draw[-{Stealth[length=2mm]},thick] (3.65,.23)--(4.25,.23);
  \draw[-{Stealth[length=2mm]},thick] (8.9,.23)--(9.5,.23);
\end{tikzpicture}
\caption{Bounding polygons in the diagonal compression.  The shear turns
lines of direction $(1,1)$ into vertical lines.  Moving each vertical slice
to start at height zero produces $P$; the marked slices correspond under
these operations.  These are continuous bounds: actual row exponents are
replaced by possible vanishing orders, not translated point by point.
All three polygons have area $H/2$.  Axis scales are chosen for legibility.}
\label{fig:node-compression}
\end{figure}%
}

\newcommand{\ReshapingCompressionFigure}{%
\begin{figure}[H]
\centering
\begin{tikzpicture}[x=.88cm,y=.88cm,font=\small,
  polyshape/.style={draw=teal!65!black,fill=teal!10,line width=.8pt},
  slice/.style={draw=orange!85!black,line width=1.2pt}]
  \begin{scope}
    \node at (1.5,3.25) {in the basis $(u,u')$};
    \path[polyshape] (1,0)--(3,.9)--(.4,2.4)--cycle;
    \draw[densely dashed,gray] (-.1,0)--(1,0);
    \draw[densely dashed,gray] (-.1,.9)--(.775,.9);
    \draw[densely dashed,gray] (-.1,2.4)--(.4,2.4);
    \node[left] at (-.1,0) {$q_0$};
    \node[left] at (-.1,.9) {$q_1$};
    \node[left] at (-.1,2.4) {$q_2$};
    \draw[slice] (.775,.9)--(3,.9);
    \node[above] at (1.9,.9) {$w$};
    \node at (1.5,-.62) {height range $R$};
  \end{scope}
  \begin{scope}[xshift=4.65cm]
    \node at (1.35,3.25) {horizontal compression};
    \draw[->,gray!65] (0,-.12)--(0,2.75);
    \path[polyshape] (0,0)--(2.225,.9)--(0,2.4)--cycle;
    \node[left] at (0,0) {$q_0$};
    \node[left] at (0,.9) {$q_1$};
    \node[left] at (0,2.4) {$q_2$};
    \draw[slice] (0,.9)--(2.225,.9);
    \node[right] at (2.225,.9) {$(w,q_1)$};
    \draw[gray!70,densely dashed] (.85,.344)--(.85,1.827);
    \node at (1.35,-.62) {maximum slice length $w$};
  \end{scope}
  \begin{scope}[xshift=9.3cm]
    \node at (1.4,3.25) {vertical compression};
    \draw[->,gray!65] (-.12,0)--(2.65,0);
    \draw[->,gray!65] (0,-.12)--(0,2.75);
    \path[polyshape] (0,0)--(2.225,0)--(0,2.4)--cycle;
    \node[below left] at (0,0) {$0$};
    \node[below] at (2.225,0) {$w$};
    \node[left] at (0,2.4) {$R$};
    \draw[slice] (.85,0)--(.85,1.483);
    \node at (1.4,-.62) {$\Delta(w,R)$};
  \end{scope}
  \draw[-{Stealth[length=2mm]},thick] (3.55,1.65)--(4.15,1.65);
  \draw[-{Stealth[length=2mm]},thick] (8.75,1.65)--(9.35,1.65);
\end{tikzpicture}
\caption{The primitive direction is chosen so that the transverse height
range is $R=H/w$.  Since the area is $H/2$, the longest horizontal slice
has length $w$.  Horizontal compression aligns the left endpoints of
these slices.  Vertical compression then gives the right triangle with
intercepts $w$ and $R$.  The middle and right panels mark vertical slices
of equal length; the drawing is schematic.}
\label{fig:final-compressions}
\end{figure}%
}


\begin{proof}
Fix the target $w$.  We first use a node to replace the triangular bound
$\Delta(A,B)$ by a quadrilateral of the same area.  A diagonal compression
then produces a triangle with a variable shape.  We choose that shape and a
primitive lattice direction together so that two further compressions give
exactly $\Delta(w,H/w)$.

Throughout, write $\gamma=(\gamma_3,\ldots,\gamma_n)$ for an exponent in
the unchanged coordinates.  For a label occurring in degree $k$, the pair
exponents with this label lie in
\begin{equation}\label{eq:residual-degree}
 \kappa\Delta(A,B),\qquad
 \kappa=k-\sum_{i=3}^n\frac{\gamma_i}{a_i}\ge0.
\end{equation}
For $n=2$ the label is empty and $\kappa=k$.  Each operation below acts on
these two-coordinate spaces, carrying the common monomial with exponent
$\gamma$ along unchanged.  In particular, a planar bound by $\kappa K$
will hold simultaneously for every label.  If $\kappa=0$, the only possible
pair exponent is zero and it stays zero throughout; hence the geometric
calculations may be made with $\kappa>0$.

\paragraph{The node and its exponent bound.}
In pair variables $x,y$, consider
\begin{equation}\label{eq:node-polynomial}
 g(x,y)=y^2-x^2(1+x).
\end{equation}
This polynomial is irreducible, since $1+x$ is not a square in $\C(x)$.
Its curve has the parametrization
\begin{equation}\label{eq:node-parametrization}
 x=v^2-1,\qquad y=v(v^2-1),
\end{equation}
which covers the part where $x\ne0$.  Near the origin choose the analytic
square root with value $1$ and set
\begin{equation}\label{eq:node-coordinates}
 \xi=y+x\sqrt{1+x},\qquad
 \eta=y-x\sqrt{1+x}.
\end{equation}
These are local coordinates and $g=\xi\eta$.

For now fix a real parameter $t>0$, and order the Taylor exponents $(\alpha,\beta)$
in $\xi,\eta$ by lexicographic comparison of
$(\alpha+t\beta,\beta)$, and denote the least exponent by $\nu_t$.
This is a multiplicative well-order: only finitely many nonnegative integer
exponents have weight below any given bound.  The node gives two controls
on these exponents: its two branches bound the least weight after factors
of $g$ have been removed, while each extracted factor translates the least
exponent by $(1,1)$.

Let $f\ne0$ be supported in $\kappa\Delta(A,B)$, and write $f=g^j f_0$,
where $j\ge0$ and $g$ does not divide $f_0$.
The weighted polynomial degree with weights $(1/A,1/B)$ is additive on
nonzero products, since highest-weight forms multiply nontrivially.  As
$g$ has weighted degree $d$, the weighted degree of $f_0$ is at most
$\kappa-jd$, which is nonnegative.  Under
\eqref{eq:node-parametrization}, its restriction is a nonzero polynomial
in $v$: otherwise $f_0$ would vanish on the irreducible curve $g=0$.
Define
\begin{equation}\label{eq:normalization-degree-constant}
 W=\max(2A,3B)=dH.
\end{equation}
The monomial inequality
\[
 2i+3i'\le W\left(\frac{i}{A}+\frac{i'}{B}\right)
\]
shows that the restriction polynomial has degree at most
\begin{equation}\label{eq:restriction-degree}
 (\kappa-jd)W.
\end{equation}

Let $\mu$ be the weight of $\nu_t(f_0)$.  At $v=1$, the branch is
$\eta=0$ and $\xi$ vanishes simply.  At $v=-1$, the branch is $\xi=0$
and $\eta$ vanishes simply.  Indeed, the chosen square root becomes $v$
near $1$ and $-v$ near $-1$, and the nonzero branch coordinate in either
case is $2v(v^2-1)$.  Every nonzero pure $\xi$ term in $f_0$ has exponent
at least $\mu$, and every nonzero pure $\eta$ term has exponent at least
$\mu/t$.  The restriction polynomial therefore has orders at least
$\mu$ and $\mu/t$ at these two distinct points.  Comparing their sum with
\eqref{eq:restriction-degree} gives
\begin{equation}\label{eq:two-branch-bound}
 \left(1+\frac1t\right)\mu\le(\kappa-jd)W.
\end{equation}
To express this bound as a triangle, define its intercepts
\begin{equation}\label{eq:node-intercepts}
 a=\frac{Wt}{1+t},\qquad b=\frac{W}{1+t}.
\end{equation}
For $\nu_t(f_0)=(\alpha,\beta)$, inequality~\eqref{eq:two-branch-bound}
is exactly $\alpha/a+\beta/b\le\kappa-jd$.  Hence
$\nu_t(f_0)\in(\kappa-jd)\Delta(a,b)$.

The extracted factors account for the rest of the degree bound.  Set
$c=1/d$, so their contribution is
$\nu_t(g^j)=(j,j)=jd(c,c)$.  The coefficients $\kappa-jd$ and $jd$
sum to $\kappa$, and therefore
\begin{equation}\label{eq:node-envelope}
 \begin{split}
 \nu_t(f)&\in(\kappa-jd)\Delta(a,b)+jd\{(c,c)\}\subset\kappa Q,\\
 Q&=\conv\{(0,0),(a,0),(c,c),(0,b)\}.
 \end{split}
\end{equation}
Thus the axis triangle comes from the two branch orders, and the diagonal
point comes from powers of the node equation.

From this point on, restrict the free parameter to $t>T$, where
\begin{equation}\label{eq:node-parameters}
 T=d^2H=\max(9B/A,4A/B)\ge6.
\end{equation}
The inequality
\[
 \frac ca+\frac cb=\frac{(1+t)^2}{Tt}>1
\]
shows that $Q$ is a quadrilateral with the vertices in the stated order.
Its area is
\begin{equation}\label{eq:node-area}
 \operatorname{area}(Q)=\frac{c(a+b)}2=\frac H2.
\end{equation}
The identity $W=dH$ is what makes this area equal to that of the input
triangle $\Delta(A,B)$.  The parameter $t$ remains free to change the
shape of $Q$ for the later lattice-direction choice.
For each label, replace its pair support by the possible values of
$\nu_t$ on its polynomial span.  The product of the spaces in degrees
$k,l$ with labels $\gamma,\gamma'$ lies in the space of degree $k+l$
with label $\gamma+\gamma'$, by the original support inclusion.
Lemmas~\ref{lem:least-terms}
and~\ref{lem:row-valuation} preserve cardinality and the graded inclusion,
and \eqref{eq:node-envelope} supplies the bound $\kappa Q$.

We also need the jet implication for this change.  At each fixed degree,
choose a basis with distinct least exponents in each nonempty label space.
One sufficiently small $e>0$ makes weights $(1,t+e)$ select all these
least terms strictly.  To see this, choose $M$ larger than all the selected
$(1,t)$-weights.  Among the finitely many exponents of weight at most $M$,
a small positive perturbation preserves strict comparisons and realizes
the tie-break by the second exponent.  Make $e$ smaller still so that all
selected weights remain below $M$; exponents originally above $M$ cannot
interfere, because a positive perturbation only increases their weights.
Lemmas~\ref{lem:weighted-limit}
and~\ref{lem:rank-transfer} apply under
$(\xi,\eta)=(sX,s^{t+e}Y)$, with the other variables unchanged.
Their inverse images are torus points.  In fact, for each of the finitely
many test points with $X,Y\ne0$, one has $\eta/\xi\to0$, while
\[
 x=\frac{\xi-\eta}{2}+O\bigl((\xi-\eta)^2\bigr),
 \qquad y=\frac{\xi+\eta}{2}.
\]
Thus $x,y\ne0$ for all sufficiently small positive $s$.  The local
coordinate map is injective on a common neighborhood, so it also preserves
distinctness of the test points.

\paragraph{From the quadrilateral to a triangle.}
Apply the integral shear
\[
 (\alpha,\beta)\longmapsto(\alpha-\beta,\beta)
\]
and then compress in the second coordinate.  The transformed quadrilateral
has vertices $(0,0),(a,0),(0,c),(-b,b)$.  For first coordinate $q$ in
$[-b,0]$, its lower boundary is $-q$ and its upper boundary is
$c+(c-b)q/b$; their difference is $c(1+q/b)$.  For $q\in[0,a]$, the
vertical slice length is $c(1-q/a)$.  Compression therefore gives the
triangle
\begin{equation}\label{eq:compressed-triangle}
 P=\conv\{(-b,0),(a,0),(0,c)\},\qquad
 \operatorname{area}(P)=\frac H2.
\end{equation}
Figure~\ref{fig:node-compression} displays these slice lengths geometrically.
Lemmas~\ref{lem:unimodular} and~\ref{lem:compression} show that the exponent
sets remain graded, retain their cardinalities, admit backward jet
transfer, and are bounded, label by label, by $\kappa P$.

\ReshapingNodeFigure

The area is now correct, but we have not yet obtained the prescribed
intercept $w$.  We choose a primitive integral direction whose transverse
range is exactly $H/w$.  In the resulting coordinates, the longest
horizontal slice will have length $w$; compression will turn that length
into the horizontal intercept.

\paragraph{Choosing the primitive direction.}
Put $R=H/w$ and $a_0=WT/(1+T)$.  As $t$ ranges over $(T,\infty)$,
the intercept $a$ ranges over $(a_0,W)$.  Choose $j$ to be the largest
power of two not exceeding
\[
 \frac{R}{2W}=\frac{1}{2dw}.
\]
Then
\begin{equation}\label{eq:primitive-j}
 j>\frac{1}{4dw}>4,
 \qquad R-jW\ge\frac R2>0.
\end{equation}
The open interval
\begin{equation}\label{eq:primitive-interval}
 \bigl(d(R-jW),\ d(R-ja_0)\bigr)
\end{equation}
has length
\[
 \frac{djW}{1+T}=\frac{jT}{1+T}>2.
\]
It therefore contains an odd integer $i$.  This is the point of choosing
$j$ to be a power of two: the vector $u=(i,j)$ is automatically primitive.
Choose the unique $t>T$ for which
\begin{equation}\label{eq:tuned-intercept}
 a=\frac{R-i/d}{j}.
\end{equation}
The interval \eqref{eq:primitive-interval} ensures that this value lies in
$(a_0,W)$.

At the three vertices $(-b,0),(a,0),(0,c)$ of $P$, the functional
$z\mapsto\det(u,z)$ takes values $jb,-ja,i/d$, respectively.  They are
distinct and ordered as
\begin{equation}\label{eq:vertex-heights}
 -ja<jb<i/d,
 \qquad i/d-jb=R-jW>0.
\end{equation}
Their range is exactly $i/d+ja=R$.  This choice used only inequalities
between real numbers; in particular, no rationality of the intercepts was
required.  We now fix this $t$ and this primitive vector for all degrees.

\paragraph{The two final compressions.}
Complete $u$ to an integral basis $(u,u')$ with determinant $1$, and use
the coefficients in that basis as new exponent coordinates.  The second
coordinate of a vector $z$ is $\det(u,z)$.  Thus the transformed triangle
still has area $H/2$, and its three vertex heights, say
$q_0<q_1<q_2$, satisfy $q_2-q_0=R$.

Its horizontal slice length is zero at $q_0$, increases linearly up to
$q_1$, and decreases linearly to zero at $q_2$.  If the maximum slice
length is $h_{\max}$, integration of this tent gives area
$Rh_{\max}/2$.  Hence $h_{\max}=H/R=w$.  Compressing in the first
coordinate therefore gives
\begin{equation}\label{eq:tent-triangle}
 \conv\{(0,q_0),(w,q_1),(0,q_2)\}.
\end{equation}
Its vertical slice at first coordinate $h\in[0,w]$ has length
$R(1-h/w)$.  A final compression in the second coordinate produces
$\Delta(w,R)$, as illustrated in Figure~\ref{fig:final-compressions}.

\ReshapingCompressionFigure

All these changes satisfy Lemmas~\ref{lem:unimodular}
and~\ref{lem:compression}.  They commute with scaling of the bounding
shapes, so for the label $\gamma$ the final pair exponents lie in
$\kappa\Delta(w,R)$.  Substituting \eqref{eq:residual-degree} gives
\eqref{eq:reshape-bound}.  Cardinality, the graded inclusion, and the
backward jet implication hold at every stage, proving all the assertions.

For clarity, the sets in every degree are defined by the fixed exact order
$\nu_t$, the fixed integral maps, and exact orders of vanishing at $1$.
The perturbation $e$ and the small scaling parameters only witness
jet-rank implications at a specified degree.  They may depend on that
degree and on the test tuple without changing the sets or their graded
inclusion.
\end{proof}
